% Preamble của báo cáo. Nạp bởi report.tex, không tự biên dịch được.
%
% BẮT BUỘC dùng XeLaTeX, không phải pdfLaTeX. Lý do nằm ở decision record trong
% docs/plans/active/2026-10-04-bao-cao-latex-plan.md: tiếng Việt cần Unicode đầy đủ và một
% font có đủ dấu tổ hợp, mà đường encoding 8-bit của pdfLaTeX vỡ ngay khi gặp một ký tự
% ngoài bảng mã.
%
% Font hệ thống mà bản build cần: Times New Roman (chữ) và Consolas (mã). Cả hai có sẵn trên
% Windows. Trên máy không có chúng, XeLaTeX sẽ dừng với lỗi "cannot find font" -- khi đó đổi
% hai dòng \setmainfont và \setmonofont bên dưới, đừng đổi gì khác.

\usepackage{fontspec}
\usepackage{polyglossia}
\setmainlanguage{vietnamese}
\setmainfont{Times New Roman}
\setmonofont{Consolas}[Scale=0.88]

\usepackage[a4paper,inner=3cm,outer=2.5cm,top=2.5cm,bottom=2.5cm]{geometry}

% polyglossia khai báo tiếng Việt là `nohyphenation`, và đúng như vậy: tiếng Việt không ngắt
% âm tiết bằng gạch nối. Hệ quả là TeX mất hẳn công cụ chính để chỉnh một dòng quá dài, nên
% văn bản tiếng Việt không nới gì sẽ đầy overfull hbox -- chữ tràn ra ngoài lề phải. Hai dòng
% dưới trả lại cho nó chỗ co giãn ở khoảng trắng giữa từ.
\setlength{\emergencystretch}{3em}
\tolerance=2000

\usepackage{graphicx}
\graphicspath{{figures/}}
% Trang ngang cho các sơ đồ dày. Kích thước bên trong \begin{landscape}, đo bằng \typeout
% chứ không suy ra: \textwidth giữ nguyên 441pt, \textheight bị gán **thành** 441pt, và thứ
% rộng 702,8pt là \linewidth. Nên bề rộng hình khổ ngang là `width=\linewidth`, còn trần
% chiều cao là `height=\textheight`. Dùng \textheight làm bề rộng -- cái tên nghe có vẻ
% đúng -- cho ra hình hẹp hơn đáng có 37%.
\usepackage{pdflscape}
\usepackage{etoolbox}           % \ifdefempty, dùng cho các dòng bìa có thể để trống
\usepackage[font=small,labelfont=bf]{caption}

\usepackage{booktabs}
\usepackage{array}
\usepackage{longtable}          % bảng 13 bảng dữ liệu ở chương 7 vắt qua trang
\usepackage{enumitem}
\setlist{nosep,leftmargin=*}

\usepackage[unicode,hidelinks]{hyperref}
\hypersetup{
  pdftitle={Kriky — Trợ lí ra đề và chữa bài},
  pdflang={vi},
}

% Một cột hẹp cố định, căn trên, dùng cho cột đầu của các bảng tra ở chương 5 và 7.
\newcolumntype{L}[1]{>{\raggedright\arraybackslash}p{#1}}
